% Cote 119: offprints, reports, courses, CV and letters (typed and printed part): archive cote 119, site page 288 (PDF page 289).
% typescript. Draft compared: none.
% AI-transcribed by Claude Opus 5.5 (claude-opus-5-5, adaptive thinking, effort high), Message Batches API, batch r119_full (msgbatch_01DwZNm5Y7YsKBWyAN4sFs8w), 2026-11-13; not yet checked by eye.
\scanpage{119/288}
\textbf{COLLÈGE DE FRANCE - \ Chaire d'Algèbre et Géométrie}

Jean-Pierre SERRE

\textbf{Résumé des cours de l'année 1990-1991}

Le cours a été consacré au même sujet que celui de 1962-1963 : la \uline{cohomologie}
\uline{galoisienne}. Il a surtout insisté sur les nombreux problèmes que posent les
groupes semi-simples lorsque l'on ne fait pas d'hypothèse restrictive sur
le corps de base.

\textbf{§1. Notations}

- \ $k$ \ est un corps commutatif, supposé de caractéristique $\neq 2$, pour simplifier;
- $k_s$ \ est une clôture séparable de \ $k$;
- \ $\mathrm{Gal}(k_s/k)$ \ est le groupe de Galois de \ $k_s/k$; c'est un groupe profini.
Si \ $G$ \ est un groupe algébrique sur \ $k$, on note \ $H^1(k, G)$ \ le premier
ensemble de cohomologie de \ $\mathrm{Gal}(k_s/k)$ \ à valeurs dans \ $G(k_s)$, cf. \uline{Cohomologie}
\uline{Galoisienne}, LN 5, p.I-56. C'est un ensemble pointé.

Si \ $A$ \ est un \ $\mathrm{Gal}(k_s/k)$-module, on définit pour tout \ $n \geqslant 0$ \ des groupes de
cohomologie
\[ H^n(k, A) \ = \ H^n(\mathrm{Gal}(k_s/k), A), \]
cf. LN 5, p.I-9.
Par exemple, si \ $A = Z/2Z$, on a
\[ H^1(k, Z/2Z) \ = \ k^*/k^{*2} \]
et
\quad $H^2(k, Z/2Z) = \mathrm{Br}_2(k)$ \ (noyau de la multiplication par \ $2$ \ dans le groupe
de Brauer de \ $k$).

L'un des thèmes du cours a été d'expliciter les relations qui existent
(ou qui pourraient exister) entre l'ensemble \ $H^1(k, G)$, pour \ $G$ \ semi-simple,
et les groupes \ $H^n(k, A)$ \ pour \ $A = Z/2Z$ (ou \ $Z/3Z$, ou tout autre `` petit''
module sur \ $\mathrm{Gal}(k_s/k)$).

\mspage{\hand{288}}
